\subsection{Relational Projection and Admissibility Principles}
  \label{subsec:relational-projection}

  Admissible physical descriptions arise through a projection from the
  relational substrate~$\chi$ onto effective observable configurations.
  This projection does not operate directly on individual configurations
  but on their relational structure.

  \paragraph{Relational operator.}
    Let $\mathcal{H}_\chi$ denote the configuration space of admissible
    $\chi$ states.
    We introduce a relational operator
    \begin{equation}
      L_\chi : \mathcal{H}_\chi \rightarrow \mathcal{H}_\chi ,
    \end{equation}
    defined purely in terms of correlations within the $\chi$ substrate.
    In discrete realizations (Appendix~D) $L_\chi$ reduces to a graph
    Laplacian.
    In the continuum limit it defines an effective self-adjoint operator
    encoding relational connectivity.

  \paragraph{Spectral structure.}
    The operator $L_\chi$ admits a spectral decomposition
    \begin{equation}
      L_\chi \psi_n = \lambda_n \psi_n ,
    \end{equation}
    where $\{\lambda_n\}$ are non-negative eigenvalues and $\{\psi_n\}$ the
    associated eigenmodes.
    These modes provide a natural basis for describing relational
    configurations of the substrate.

  \paragraph{Admissibility principle.}
    Only a restricted subset of relational configurations can give rise
    to stable effective descriptions.
    This restriction can be expressed as a spectral filtering acting on
    the relational operator $L_\chi$.
    Formally one may represent this selection through an admissibility
    filter
    \begin{equation}
      F_{\lambda_*} : \mathcal{H}_\chi \rightarrow \mathcal{H}_{\mathrm{eff}},
    \end{equation}
    which associates relational configurations with admissible effective
    descriptions.
    The scale $\lambda_*$ represents a spectral threshold separating
    projectively admissible from non-admissible relational modes.

    Modes with $\lambda_n \le \lambda_*$ contribute directly to effective
    descriptions, while higher modes remain encoded in the fibre of the
    projection and appear only indirectly through fluctuations or
    renormalized parameters.
    Such modes constitute infra-projectable invariants of the relational
    description.

  \paragraph{Non-injectivity and equivalence classes.}
    The admissibility filter is generally non-injective.
    Distinct relational configurations may share identical spectral
    content below the admissibility threshold and therefore correspond
    to the same effective observable description.
    Effective states therefore arise as equivalence classes of relational
    configurations compatible with a common observable structure.

  \paragraph{Programme of spectral admissibility.}
    The structural consequences of this admissibility principle are
    developed in detail in Chapter~\ref{sec:spectral-admissibility}.
    In particular, the bounded-relaxation dynamics of the $\chi$ substrate
    imposes quantitative constraints on admissible spectral modes through
    an admissibility envelope that limits their effective amplitudes.

    This constraint induces a hierarchy among spectral sectors and leads
    to a classification of relational structures capable of supporting
    stable projected descriptions.
    In later sections this mechanism will be shown to select specific
    finite symmetry groups compatible with admissible relational dynamics.
